%% text-growth.tex -- the growth stage (used by florilegium-grow.tex only)
\subsubsection{The growth stage}

This document is the answer to the exercise. Its margin is not the authored composition of
the previous section: nothing in it was placed by hand. A Lua program decides where every
stem, leaf and flower goes and how large it is, and then calls the same primitives to paint
them.

\begin{construction}[Stems]
The root is placed near the bottom-left corner of the page. Three stems rise from it; the
$x$ coordinate of each is a smooth function of the height,
\[
  x(y)=c+A\,n\bigl(y/\lambda+\varphi\bigr),
  \qquad A=0.2\,\text{(column width)},\ \lambda\in[6,9]\,\text{cm},
\]
with $n$ a two-octave gradient noise and $\varphi$ a phase drawn per stem, blended into the
root over the first four centimeters. Different phases make the stems sway apart and cross
now and then, as real stems do. The first reaches the top of the page and ends in the
largest flower; the second stops at $60$ to $72\%$ of the height and ends in a lily; the
third at $36$ to $48\%$, in a flower. A footer stem leaves the root along the bottom of the
page. Leaves are hung at the nodes, every $1.4$ to $2.2\,$cm, alternately left and right,
$35^\circ$ to $60^\circ$ off the stem, $2.4\,$cm long at the bottom of the page and $1.3\,$cm at
the top, and only where there is room.
\end{construction}

\begin{definition}
A \emph{proxy} for a motif is a disc of radius $r(s)$, where $s$ is the motif's scale, that
approximates the area it occupies. A point $p$ of the free region $R$ has \emph{free
radius}
\[
  f(p)=\min\Bigl(\,d(p,\partial R),\ \min_i\bigl(\lVert p-c_i\rVert-\kappa\,r_i\bigr)\Bigr),
\]
the radius of the largest disc about $p$ that stays inside $R$ and overlaps no proxy
$(c_i,r_i)$ by more than the fraction $1-\kappa$. With $\kappa=0.85$, a leaf may tuck under
the edge of a flower; stems are proxies too, as chains of small discs.
\end{definition}

\begin{construction}[Growth]
The region $R$ is the margin column, a band along the bottom of the page, and on the last
page also the rectangle between the end of the text and the footer. On a grid of candidates
spaced $2\,$mm apart, each point keeps $f(p)$ and its distance to the nearest stem. Then,
until the largest free radius falls below $3.6\,$mm:
\begin{enumerate}
  \item take the point $p$ that maximizes $f(p)$ among those within $3\,$cm of a stem, with
        a random factor in $[0.85,1.15]$ so that ties break differently on every seed;
  \item choose a motif whose smallest proxy fits, by weight; a motif that has already
        appeared $k$ times has its weight divided by $1+k/c$, and only flowers, lilies,
        umbels and bells have a hard cap. Set its scale so that $r(s)\approx0.95\,f(p)$;
  \item give it an axis pointing away from the nearest stem point, leaning upward for the
        kinds that reach for the light, and grow a branch to it from that point, as in the
        authored composition. The branch becomes a stem for later branches.
\end{enumerate}
Because the largest hole is always taken first, flowers and lilies are placed early,
sprays and leaves fill the middle, and a last pass puts pairs of leaflets, buds and small
clusters of florets on the stems wherever a hole of more than $1.7\,$mm is left next to
one; a hole far from any stem may get one of six copper accents, and otherwise stays empty.
That ordering is the whole reason the result looks \emph{composed} rather than scattered.
\end{construction}

\begin{remark}[Adapting to the text]
The free region is computed from the page layout, so changing the margins moves the
garden. The end of the text is measured where it happens: \verb|\florilegiumtextend|
records \verb|\pagetotal| and \verb|\pagedepth|, the height of the material already on the
current page, and the background of that page is drawn after it. On the last page, the
garden therefore spills into the space below the words in a single compilation: branches
arch out of the main stems into that space and each ends in a flower.
\end{remark}

\begin{remark}[Visitors, marks and seasons]
Butterflies and bees take holes like any motif but need no branch; a snail is put on a
main stem and, on about one page in five, a bird on a twig of its own. Their decisions come
from a random stream of their own, so a page without them is the page it would have been.
A mark in the text, left by \verb|\florilegiumbloom| or by a section heading, is found in
the page box at shipout and gets its flower before any leaf is hung. An age between $0$
and $1$ reweights the motifs from buds in spring to seed heads in autumn.
\end{remark}

\begin{remark}[Cost]
The layout is cheap: the grid has a few thousand points, a proxy only updates the points
near it, and the whole loop takes a fraction of a second. This document of four pages
compiles in about five seconds.
\end{remark}
